\section{Introduction}\label{sec:intro}

Optimization models in process systems engineering, energy networks and
optimal control are often built from local pieces. A flowsheet links unit
models along streams, a multiperiod model links consecutive time stages,
and a network model links quantities at adjacent nodes. The resulting
objective is a sum of terms, each depending on a few variables, and the
graph that records which variables share a term often has small treewidth.
For variables with finitely many values, small treewidth is exactly what
dynamic programming needs: bucket elimination or junction-tree message
passing finds a global optimum in time exponential only in the width of a
tree decomposition \cite{Bertele1972,Dechter1999,KollerFriedman2009}.

Mixed-integer nonlinear models have continuous variables and nonconvex
terms, and there the situation is less clear. Global solvers based on
spatial branch-and-bound compute bounds from convex relaxations of
individual terms and branch on variable domains
\cite{TawarmalaniSahinidis2005,Belotti2009,Misener2014}; they do not use a
tree decomposition to organize the search. Discretizing the continuous
variables and running a dynamic program gives neither a valid bound nor a
running time that scales well with the requested accuracy: a uniform grid
of mesh $h$ has $\Theta(1/h)$ points per coordinate, so tables grow like
$h^{-p}$ in the bag size $p$. Sparse polynomial optimization methods
exploit the same structure through moment and linear-programming
hierarchies \cite{Waki2006,Lasserre2006,BienstockMunoz2018}, with sizes
that also grow polynomially in the inverse accuracy.

Small treewidth alone does not make continuous nonconvex problems
tractable. Del Pia and Khajavirad \cite{DelPiaKhajavirad2026} recently
showed that box-constrained nonconvex quadratic programming is solvable in
strongly polynomial time when the interaction graph is a forest, but is
strongly NP-hard at treewidth two, even with small integer coefficients;
quartic minimization is strongly NP-hard on paths. This paper asks which
additional, checkable structure makes tree decompositions useful for
\emph{certified} global optimization of mixed-integer problems, and what
fails when that structure is absent.

\subsection{Main results}
We consider $\min\{F(x):x\in X\}$, where $F=\sum_af_a(x_{S_a})$ is a sum of
factors, $X$ is a box in which some coordinates are integer, and a tree
decomposition of the scopes $S_a$ with maximum bag size $p$ is given. Two
quantities govern our results: an upper bound $L$ on the curvature of $F$
along each coordinate (the diagonal of the Hessian for a quadratic), and a
quadratic growth constant $g$, defined by
$F(x)-\OPT\ge g\norm{x-x^*}^2$ for all $x\in X$, where $x^*$ is the unique
minimizer. We write $\kappa=\max\{1,L/g\}$.

\paragraph{Certified coordinate grids.}
Choose finite grids in the coordinates, and subtract from $F$ at each grid
node $v$ of coordinate $i$ the unary correction $L w^2/8$, where $w$ is the
largest grid interval ending at $v$. The minimum of the corrected objective
over the grid is a lower bound on $\min_XF$ (Proposition~\ref{prop:cellwise}).
Because the corrections are unary, the minimum is computed by ordinary
dynamic programming on the given decomposition, together with all
coordinate min-marginals. These min-marginals give a safe rule for
discarding coordinate intervals that cannot contain a point better than a
known feasible value, and the sequence of grids, min-marginals and
incumbents forms a certificate that an independent program can check
(Theorem~\ref{thm:certificate}). None of this uses convexity, uniqueness or
growth.

\paragraph{Accuracy-independent grids.}
Under quadratic growth, grids whose spacing grows linearly away from the
current corrected-grid minimizer, combined with this filtering, keep
$O(\sqrt\kappa\log n)$ nodes per coordinate at every stage, independently
of the accuracy reached (Lemma~\ref{lem:states}). A schedule of capped
trials removes the need to know $g$. For an explicit rational
mixed-integer box quadratic program this gives a certified
$2^{-q}$-approximation in $f(p,\kappa)(I+q)^{O(1)}$ bit operations, where
$I$ is the input length (Theorem~\ref{thm:approx}), and the exact rational
minimizer and optimal value in $f(p,\kappa)I^{O(1)}$ bit operations
(Theorem~\ref{thm:exact}). The algorithm is thus fixed-parameter tractable
in the pair $(p,\kappa)$. The same approximation result holds for
fixed-degree polynomial factors, with exact output when all coordinates are
integer (Corollary~\ref{cor:poly}). An exact acceptance rule based on the
denominator of a candidate's value certifies optimality of any proposed
point, and a snapping procedure turns any point whose value is within
$\min\{g_S/(64n^2R^2),1/(2\Omega^2)\}$ of $\OPT$ into an exact minimizer,
where $g_S$ is the growth constant towards the optimal \emph{set} and
$R,\Omega$ are explicit height bounds (Theorem~\ref{thm:transfer}); no
uniqueness is needed for this step.

A simple chain with box widths $2^t$ (Example~\ref{ex:chain}) illustrates
the difference from exact parametric dynamic programming: its exact Bellman
messages need at least $2^{m-1}$ quadratic pieces, while the corrected
grids certify any accuracy in time polynomial in $m$. Our implementation
certifies a gap of $10^{-3}$ for $m=64$ with at most eleven nodes per
coordinate grid (Section~\ref{sec:computation}).

\paragraph{Conditional recourse.}
The condition number uses the largest diagonal curvature, so a stiff convex
coupling term inflates $\kappa$ even though convex subproblems are easy. The
corrected-grid bound applies unchanged to any value function
$V(v)=\min_yF(v,y)$ whose inner feasible set does not depend on $v$:
partial minimization preserves coordinate curvature and growth
(Lemma~\ref{lem:valuefunction}). We use this in three ways. Private convex
quadratic blocks with a certified affine or piecewise-affine response can
be eliminated exactly, preserving the decomposition and replacing the
direct curvature by the curvature of the reduced problem; private convex
blocks without such a response become value factors evaluated by exact
convex quadratic programming. A large residual problem that is separately
concave and becomes submodular after sign changes can be removed from
discretization altogether, since each conditional value is one minimum
cut; only a small core is then searched by adaptive corrected cells. A
counterexample shows that the global correction cannot simply be replaced
by bag-local corrections of grid messages; exact or certified conditional
values are what make local corrections valid
(Section~\ref{sec:recourse}).

\paragraph{Constraints and nonunique optima.}
Independent rounding, on which the certificate rests, breaks under coupling
constraints. When the continuous coupling matrix is totally unimodular,
correlated rounding to feasible cell corners restores the certificate on
aligned uniform grids, with exact feasibility and exact output for
quadratics; the resulting algorithm is polynomial for fixed width and
$\kappa$, but not fixed-parameter tractable (Section~\ref{sec:constraints}).
For coordinatewise concave mixed-integer quadratics, the Bellman residuals
of an endpoint dynamic program give a factored description of the entire
optimal set (Section~\ref{sec:optsets}).

\paragraph{Structural limits.}
Section~\ref{sec:limits} collects what fails without each ingredient:
hardness at bounded width when $\kappa$ is unbounded and a lower bound on the
dependence on $p$; exponential exact messages (Example~\ref{ex:chain});
invalidity of local grid corrections; a family with bounded width,
bounded negative-curvature ratio and a unique minimizer for which matching
local moments does not control the error of a local relaxation; the role
of product domains; and the failure of point localization when only the
distance to an optimal \emph{set} grows.

\subsection{Relation to prior work}
\paragraph{Dynamic programming with continuous variables.}
Finite-state min-sum dynamic programming on tree decompositions and its
min-marginals are classical
\cite{Bertele1972,Dechter1999,WainwrightJaakkolaWillsky2005,KohliTorr2007}.
For continuous variables, Del Pia and Khajavirad's forest algorithm
\cite{DelPiaKhajavirad2026} and the parametric dynamic program of
Bhathena et al.\ \cite{Bhathena2026} for convex quadratics with indicators
represent messages exactly as piecewise quadratics; Example~\ref{ex:chain}
shows that such representations can be exponential at treewidth two even
with a unique, well-conditioned minimizer. Interval dynamic programming
with relaxation-based bounds has been used for optimal power flow on tree
networks \cite{Dvijotham2017}, continuous distributed constraint
optimization uses discretized or exact messages \cite{Hoang2020}, and
adaptive or variable-resolution discretization of dynamic programs and
message passing is common in practice
\cite{MunosMoore2002,LeaverFay2005,Parisot2014,Troullinos2022}. We are
not aware of a method that combines a unary correction making the
discretized dynamic program a valid lower bound, min-marginal filtering
with a checkable history, and a growth-based bound on the retained grid
that is independent of the accuracy.

\paragraph{Global optimization bounds.}
The one-dimensional interpolation error $Lw^2/8$ underlies piecewise-linear
relaxations of quadratic terms \cite{Beach2022,DongLuo2018} and
$\alpha$BB-type underestimators \cite{Adjiman1998}; optimality-based bound
tightening discards regions that cannot improve an incumbent
\cite{Belotti2009,GleixnerBertholdMuller2017}. Our corrections act on the
whole objective rather than on individual terms, which is what allows exact
dynamic programming to compute the bound, and the filtering is driven by
exact min-marginals of that bound.

\paragraph{Sparse polynomial optimization.}
Bienstock and Mu\~noz \cite{BienstockMunoz2018} give linear programming
formulations of size $O((2\pi/\varepsilon)^{\omega+1}n\log(\pi/\varepsilon))$
for polynomial problems of degree $\pi$ and intersection width $\omega$, with
scaled tolerances on feasibility and objective. Sparse moment and
sum-of-squares hierarchies exploit chordal structure
\cite{Waki2006,Lasserre2006}. These methods need no growth assumption, but
their size depends polynomially on $1/\varepsilon$; under growth our
dependence on the accuracy is polynomial in $\log(1/\varepsilon)$, and the
returned points are exactly feasible.

\paragraph{Complexity of quadratic optimization.}
Rational quadratic programs have optimal solutions of polynomial size
\cite{Vavasis1990,DelPiaDeyMolinaro2017}. Box QP with a fixed number of
negative eigenvalues is polynomial \cite{Vavasis1992}, and Del Pia
\cite{DelPia2026} approximates mixed-integer QP with fixed integer
dimension and few negative eigenvalues. Integer quadratic programming is
fixed-parameter tractable under combined structural and domain restrictions
\cite{EibenGanianKnopOrdyniak2019}. Our parameter $\kappa$ is of a different
kind: negative inertia and integer dimension may grow with $n$, and the
domain may be exponentially large. Proximity and scaling arguments give
logarithmic dependence on domain size for separable convex integer
programs \cite{HochbaumShanthikumar1990,Hunkenschroder2026}; growth plays
an analogous role for our nonconvex, nonseparable objectives.

\paragraph{Recourse, cuts and total unimodularity.}
Affine and piecewise-affine solutions of parametric convex QPs are
classical in multiparametric programming \cite{Bemporad2002}. Minimizing
pairwise submodular binary energies by minimum cuts \cite{KolmogorovZabih2004}
and the endpoint property of coordinatewise concave functions are standard,
and continuous submodular minimization is an active topic
\cite{Bach2019,BurerNatarajanWillemsen2025}. Totally unimodular systems are
box-integral \cite{HoffmanKruskal1956,ChervetGrappeRobert2021}. Our
contribution in these parts is the composition of these tools with the
corrected-grid certificate and its complexity analysis, not the tools
themselves.

\subsection{Contributions}
To the best of our knowledge, the following are new:
\begin{enumerate}[label=(\roman*)]
\item the corrected-grid lower bound for sparse mixed-integer box problems
under coordinate curvature, with min-marginal interval filtering and a
full-history certificate (Section~\ref{sec:grids});
\item the accuracy-independent grid bound under quadratic growth, and the
resulting algorithms that are fixed-parameter tractable in bag size and
condition number, for certified approximation of sparse polynomial
problems and exact solution of sparse mixed-integer box quadratic
programs, without knowledge of the growth constant
(Sections~\ref{sec:growth}--\ref{sec:exact});
\item the use of value functions, certified convex responses and
minimum-cut residual oracles inside the corrected-grid framework, with
the corresponding curvature and growth transfer
(Section~\ref{sec:recourse});
\item the explicit limiting examples of Section~\ref{sec:limits}, in
particular the expanding-box chain and the local-correction and
local-moment counterexamples;
\item an exact-arithmetic implementation whose certificates are replayed
by an independent checker (Section~\ref{sec:computation}).
\end{enumerate}
The ingredients, namely interpolation error bounds, tree dynamic
programming, min-marginals, domain filtering, rational height bounds,
parametric quadratic programming, graph cuts and total unimodularity, are
classical; the novelty lies in their combination and in the resulting
guarantees.

\paragraph{Organization.}
Section~\ref{sec:setting} fixes the model and the parameters.
Sections~\ref{sec:grids}--\ref{sec:exact} develop the corrected grids, the
growth analysis and exact output. Section~\ref{sec:recourse} treats
conditional recourse, Section~\ref{sec:constraints} coupling constraints and
Section~\ref{sec:optsets} nonunique optima. Section~\ref{sec:limits}
discusses structural limits, Section~\ref{sec:computation} the
implementation and experiments, and Section~\ref{sec:conclusion} open
problems.
